% Partie : intuitionistic-logic
% Chapitre : introduction
% Section : natural-deduction

\documentclass[../../../include/open-logic-section]{subfiles}

\begin{document}

\olfileid[fr]{int}{int}{ntd}

\olsection{Déduction naturelle}

La déduction naturelle sans la règle $\FalseCl$ est un système
standard de !!{derivation} pour la logique intuitionniste. Nous
rappelons ses règles et les motivons au moyen de l'interprétation BHK.
Pour chacune, si les prémisses ont des constructions, la conclusion
en a une aussi. La règle $\FalseInt$, qui tire n'importe quelle
formule de l'absurde, reste admise : elle n'est pas la règle
classique $\FalseCl$.

Les !!{derivation}s en déduction naturelle pouvant comporter des
hypothèses non déchargées, nous devons comprendre une !!{derivation}
de $!A$ à partir d'hypothèses !!{undischarged}s~$\Gamma$ comme une
fonction transformant des constructions de tous les $!B \in \Gamma$
en une construction de~$!A$. S'il existe une !!{derivation} de~$!A$
sans hypothèse !!{undischarged}, il existe alors une construction
de~$!A$ au sens de BHK. Dans la discussion, nous omettrons toutefois
la mention de~$\Gamma$ lorsqu'elle n'est pas nécessaire.

Une hypothèse~$!A$ constitue à elle seule une !!{derivation}
de~$!A$ à partir de l'hypothèse !!{undischarged}~$!A$. Cela
concorde avec l'interprétation BHK : la fonction identité transforme
toute construction de~$!A$ en une construction de~$!A$.

\subsection{Conjonction}

\begin{defish}
\AxiomC{$!A$}
\AxiomC{$!B$}
\RightLabel{\Intro{\land}}
\BinaryInfC{$!A \land !B$}
\DisplayProof
\hfill
\begin{tabular}{l}
\AxiomC{$!A \land !B$}
\RightLabel{\Elim{\land}}
\UnaryInfC{$!A$}
\DisplayProof
\\[3ex]
\AxiomC{$!A \land !B$}
\RightLabel{\Elim{\land}}
\UnaryInfC{$!B$}
\DisplayProof
\end{tabular}
\end{defish}

Supposons que $N_1$ et $N_2$ soient des constructions de $!A_1$
et de $!A_2$, respectivement. Nous obtenons alors une construction
de $!A_1 \land !A_2$ : le couple $\tuple{N_1, N_2}$.

Réciproquement, selon BHK, une construction de $!A_1 \land !A_2$\footnote{La source répète le premier membre de la conjonction dans cette phrase. Le second membre est rétabli conformément aux deux projections et aux types des constructions qui suivent.}
est un couple $\tuple{N_1, N_2}$. Si nous possédons un tel couple,
ses projections donnent une construction de chaque membre : $N_1$
construit~$!A_1$ et $N_2$ construit $!A_2$.

\subsection{Implication}

\begin{defish}
  \AxiomC{$\Discharge{!A}{u}$}
  \DeduceC{$!B$}
  \DischargeRule{$\Intro{\lif}$}{u}
  \UnaryInfC{$!A \lif !B$}
  \DisplayProof
\hfill
  \AxiomC{$!A \lif !B$}
  \AxiomC{$!A$}
  \RightLabel{$\Elim{\lif}$}
  \BinaryInfC{$!B$}
  \DisplayProof
\end{defish}

Si nous avons une !!{derivation} de~$!B$ à partir de
l'hypothèse !!{undischarged}~$!A$, alors nous disposons d'une
fonction~$f$ transformant les constructions de~$!A$ en
constructions de~$!B$. Cette même fonction est une construction
de $!A \lif !B$. Ainsi, si la prémisse de $\Intro\lif$ a une
construction conditionnée par une construction de~$!A$, sa
conclusion $!A \lif !B$ en a une aussi.

Inversement, supposons que $N$ construise $!A$ et que $f$
construise $!A \lif !B$. Une construction de $!A \lif !B$
est une fonction qui transforme les constructions de $!A$ en
constructions de~$!B$. Le résultat $f(N)$ est donc une
construction de~$!B$, c'est-à-dire de la conclusion de
$\Elim\lif$.

\subsection{Disjonction}

\begin{defish}
\begin{tabular}{l}
\AxiomC{$!A$}
\RightLabel{\Intro{\lor}}
\UnaryInfC{$!A \lor !B$}
\DisplayProof
\\[3ex]
\AxiomC{$!B$}
\RightLabel{\Intro{\lor}}
\UnaryInfC{$!A \lor !B$}
\DisplayProof
\end{tabular}
\hfill
\AxiomC{$!A \lor !B$}
\AxiomC{$\Discharge{!A}{n}$}
\DeduceC{$!C$}
\AxiomC{$\Discharge{!B}{n}$}
\DeduceC{$!C$}
\DischargeRule{\Elim{\lor}}{n}
\TrinaryInfC{$!C$}
\DisplayProof
\end{defish}

À partir d'une construction~$N_i$ de~$!A_i$, nous formons une
construction~$\tuple{i, N_i}$ de~$!A_1 \lor !A_2$. Inversement,
supposons donnée une construction de~$!A_1 \lor !A_2$, soit un
couple $\tuple{i, N_i}$ où $N_i$ construit~$!A_i$, ainsi que
des fonctions $f_1$ et $f_2$ transformant, respectivement, les
constructions de $!A_1$ et de $!A_2$ en constructions de~$!C$.
Alors $f_i(N_i)$ construit~$!C$, la conclusion de~$\Elim\lor$.

\subsection{Absurde}

\begin{defish}
  \AxiomC{$\lfalse$}
  \RightLabel{\FalseInt}
  \UnaryInfC{$!A$}
  \DisplayProof
\end{defish}

S'il existe une !!{derivation} de~$\lfalse$ à partir des
hypothèses !!{undischarged}s~$!B_1$, \dots, $!B_n$, alors une
fonction~$f(M_1, \dots, M_n)$ transforme des constructions de
$!B_1$, \dots, $!B_n$ en une construction de~$\lfalse$.
Comme $\lfalse$ n'a aucune construction, il ne peut y avoir
simultanément de constructions de tous les $!B_1$, \dots, $!B_n$.
Par conséquent, $f$ possède aussi la propriété conditionnelle
suivante : \emph{si} $M_1$, \dots, $M_n$ construisent, respectivement,
$!B_1$, \dots, $!B_n$, \emph{alors} $f(M_1, \dots, M_n)$ construit
également~$!A$. Cela ne postule pas l'existence d'une construction
de l'absurde.

\subsection{Règles pour $\lnot$}

Puisque $\lnot !A$ est défini comme $!A \lif \lfalse$, des règles
propres à~$\lnot$ ne sont pas nécessaires à strictement parler.
Voici néanmoins la forme qu'elles auraient :

\begin{defish}
\AxiomC{$\Discharge{!A}{n}$}
\noLine
\DeduceC{$\lfalse$}
\DischargeRule{\Intro{\lnot}}{n}
\UnaryInfC{$\lnot !A$}
\DisplayProof
\hfill
\AxiomC{$\lnot !A$}
\AxiomC{$!A$}
\RightLabel{\Elim{\lnot}}
\BinaryInfC{$\lfalse$}
\DisplayProof
\end{defish}


\subsection{Exemples de \usetoken{p}{derivation}}

\begin{enumerate}
\item $\Proves !A \lif (\lnot !A \lif \lfalse)$, c'est-à-dire $\Proves !A
  \lif ((!A \lif \lfalse) \lif \lfalse)$
  \begin{prooftree}
    \AxiomC{$\Discharge{!A}{2}$}
    \AxiomC{$\Discharge{!A \lif \lfalse}{1}$}
    \RightLabel{\Elim\lif}
    \BinaryInfC{$\lfalse$}
    \DischargeRule{\Intro\lif}{1}
    \UnaryInfC{$(!A \lif \lfalse)\lif \lfalse$}
    \DischargeRule{\Intro\lif}{2}
    \UnaryInfC{$!A \lif (!A \lif \lfalse) \lif \lfalse$}
  \end{prooftree}

\item $\Proves ((!A \land !B) \lif !C) \lif (!A \lif (!B \lif !C))$
  \begin{prooftree}
    \AxiomC{$\Discharge{(!A \land !B) \lif !C}{3}$}
    \AxiomC{$\Discharge{!A}{2}$}
    \AxiomC{$\Discharge{!B}{1}$}
    \RightLabel{\Intro\land}
    \BinaryInfC{$!A \land !B$}
    \RightLabel{\Elim\lif}
    \BinaryInfC{$!C$}
    \DischargeRule{\Intro\lif}{1}
    \UnaryInfC{$!B \lif !C$}
    \DischargeRule{\Intro\lif}{2}
    \UnaryInfC{$!A \lif (!B \lif !C)$}
    \DischargeRule{\Intro\lif}{3}
    \UnaryInfC{$((!A \land !B) \lif !C) \lif (!A \lif (!B \lif !C))$}
  \end{prooftree}

\item $\Proves \lnot(!A \land \lnot !A)$, c'est-à-dire $\Proves (!A \land (!A
  \lif \lfalse))\lif \lfalse$
  \begin{prooftree}
    \AxiomC{$\Discharge{!A \land (!A \lif \lfalse)}{1}$}
    \RightLabel{\Elim\land}
    \UnaryInfC{$!A \lif \lfalse$}
    \AxiomC{$\Discharge{!A \land (!A \lif \lfalse)}{1}$}
    \RightLabel{\Elim\land}
    \UnaryInfC{$!A$}
    \RightLabel{\Elim\lif}
    \BinaryInfC{$\lfalse$}
    \DischargeRule{\Intro\lif}{1}
    \UnaryInfC{$(!A \land (!A \lif \lfalse))\lif \lfalse$}
  \end{prooftree}

\Needspace{18\baselineskip}
\item $\Proves \lnot\lnot(!A \lor \lnot !A)$, c'est-à-dire $\Proves ((!A \lor
  (!A \lif \lfalse))\lif \lfalse)\lif \lfalse$
  \begin{prooftree}\hskip-3em
    \AxiomC{$\Discharge{(!A \lor (!A \lif \lfalse))\lif \lfalse}{2}$}
    \AxiomC{$\Discharge{(!A \lor (!A \lif \lfalse))\lif \lfalse}{2}$}
    \AxiomC{$\Discharge{!A}{1}$}
    \RightLabel{\Intro\lor}
    \UnaryInfC{$!A \lor (!A \lif \lfalse)$}
    \RightLabel{\Elim\lif}
    \BinaryInfC{$\lfalse$}
    \DischargeRule{\Intro\lif}{1}
    \UnaryInfC{$!A \lif \lfalse$}
    \RightLabel{\Intro\lor}
    \UnaryInfC{$!A \lor (!A \lif \lfalse)$}
    \RightLabel{\Elim\lif}
    \insertBetweenHyps{\hskip -4em}
    \BinaryInfC{$\lfalse$}
    \DischargeRule{\Intro\lif}{2}
    \UnaryInfC{$((!A \lor (!A \lif \lfalse))\lif \lfalse)\lif \lfalse$}
  \end{prooftree}
\end{enumerate}

\begin{prop}
  Si $\Gamma \Proves !A$ en déduction naturelle intuitionniste,
  alors $\Gamma \Proves !A$ en logique classique. En particulier,
  tout théorème intuitionniste $!A$ est aussi un théorème classique.
\end{prop}

\begin{proof}
  Chacune des règles de la déduction naturelle intuitionniste est
  aussi une règle de la déduction naturelle classique. Toute
  !!{derivation} intuitionniste est donc également une
  !!{derivation} classique.
\end{proof}

\begin{prob}
  Donnez des !!{derivation}s intuitionnistes des !!{formula}s suivantes :
  \begin{enumerate}
    \item $(\lnot !A \lor !B) \lif (!A \lif !B)$
    \item $\lnot\lnot\lnot !A \lif \lnot !A$
    \item $\lnot\lnot (!A \land !B) \liff (\lnot\lnot !A\land \lnot \lnot !B)$
    \item $\lnot (!A \lor !B) \liff (\lnot !A \land \lnot !B)$
    \item $(\lnot !A \lor \lnot !B) \lif \lnot (!A \land !B)$
    \item $\lnot\lnot ( !A \land !B) \lif  (\lnot\lnot !A \lor
    \lnot\lnot !B)$
  \end{enumerate}
\end{prob}

\end{document}
